% !TeX root = ../../Book1.tex
% 纲要标题：## 第14章　域扩张、代数性与超越性
% 纲要位置：计划纲要.md，第 500 行。

\chapter{域扩张、代数性与超越性}
\label{b1:ch:14}
\BookChapterNumber{b1:ch:14}

\begin{chapterabstract}
添入代数元与添入超越元会产生不同性质的域扩张。本章用最小多项式和向量空间维数研究代数扩张，证明次数的塔式公式；再以代数独立性、超越基与函数域描述自由参数的数量。
\end{chapterabstract}

\begin{learninggoals}
\begin{itemize}
\item 用最小多项式辨认单扩张，从不可约多项式构造含根的扩张，并用塔式公式计算次数。
\item 区分有限、有限生成与代数扩张，说明有限个代数元为何生成有限扩张。
\item 使用交换引理解释超越次数为何与超越基的选择无关。
\item 从具体函数域的生成元和关系中辨认独立参数及有限代数部分。
\end{itemize}
\end{learninggoals}

从 \(\mathbb Q\) 添入 \(\sqrt2\) 后，\((1+\sqrt2)^{-1}=\sqrt2-1\)，一切四则运算都可以整理成 \(a+b\sqrt2\) 的形式，其中 \(a,b\in\mathbb Q\)。若添入的是一个超越元 \(t\)，其各次幂却不能这样用有限个固定元素表示。两种添元过程的区别，不能仅从“加入一个新符号”看出，而要考察新元满足怎样的多项式关系，以及所得域作为原域上的向量空间有多大。扩张次数和代数独立性将分别量出这些信息。

\ProseParagraphBreak

本章需要域上多项式的带余除法、不可约性、分式域，以及线性代数中的基和维数。\cref{b1:prop:field-poly-division}及\cref{b1:thm:gauss-descent}提供主要的多项式工具。一般超越基的存在性使用 Zorn 引理，有限生成扩张中的参数选择则只用有限次选取与交换。代数闭包的存在性将在下一章证明。

\ProseParagraphBreak
有限 Galois 理论主要依赖\secref{b1:sec:1401}和\secref{b1:sec:1402}的代数扩张与次数计算。\secref{b1:sec:1403}和\secref{b1:sec:1404}的超越基、超越次数与函数域不属于下一章的前提；它们用于后续独立参数的一般方程与代数几何问题。

\ProseParagraphBreak
域扩张与代数独立性的另一种组织方式，可见 Roman 的《Field Theory》\cite[Chapter 2; Sections 4.2--4.3]{Roman2006FieldTheory}。若要补习这里使用的基、维数与线性映射，可读席南华的《基础代数》第二卷\cite[第1--2章]{Xi2018BasicAlgebraII}。对函数域的几何意义感兴趣的读者，还可在读完本章后参阅 Shafarevich 的《Basic Algebraic Geometry》第一卷\cite[Chapter 1, Sections 1.3, 3.1--3.3]{Shafarevich2013BasicAGI}，了解有理函数与函数域的关系；这属于进一步的代数几何学习。

% !TeX root = ../../Book1.tex
% 纲要标题：### 14.1 域扩张及其次数
% 纲要位置：计划纲要.md，第 502 行。

\section{域扩张及其次数}
\label{b1:sec:1401}


% 纲要内容（仅作写作计划，尚非教材正文）：
%
% * 域扩张、素域与特征
% * 有限扩张与塔式公式
% * 单扩张及最小多项式
%

把多项式的一个根添入原来的域，并非只多出一个数：加、减、乘、除会随之产生更多元素。扩张次数用线性代数的语言衡量这种变化。这里的域均为交换域，子域与原域具有同一个单位；基与维数沿用线性代数中的通常约定。

\subsection{域扩张、素域与特征}
\label{b1:subsec:1401:01}

\begin{definition}\label{b1:def:field-extension}
设 \(K,L\) 为域。如果 \(K\subseteq L\)，并且 \(K\) 的加法与乘法都是 \(L\) 中相应运算的限制，则称 \(L\) 为 \(K\) 的一个\term[yu-kuozhang]{域扩张}[field extension]，记为 \(L/K\)。给定域扩张 \(L/K\)，如果子域 \(E\subseteq L\) 满足 \(K\subseteq E\)，则称 \(E\) 为 \(L/K\) 的\term[zhongjian-yu]{中间域}[intermediate field]。给定子集 \(S\subseteq L\)，域 \(L\) 中包含 \(K\cup S\) 的最小子域记为 \(K(S)\)；它等于所有包含 \(K\cup S\) 的子域之交。
\end{definition}
\symidx{\(L/K\)：域扩张}
\symidx{\(K(S)\)：由 \(K\) 与 \(S\) 生成的子域}
记号 \(L/K\) 不表示商集。若原始数据是域同态 \(\iota:K\rightarrow L\)，则它必为单射：若非零 \(a\) 在核中，则 \(\iota(1)=\iota(a)\iota(a^{-1})=0\)，与 \(\iota(1)=1\neq0\) 矛盾。因而可把 \(K\) 与像 \(\iota(K)\) 认同，再使用包含的写法；认同时仍须记住具体嵌入。

\begin{proposition}\label{b1:prop:prime-field}
每个域 \(K\) 的特征为零或一个素数 \(p\)。它包含唯一的最小子域，称为\term[su-yu]{素域}[prime field]；特征零时同构于 \(\mathbb Q\)，特征 \(p\) 时同构于 \(\mathbb F_p=\mathbb Z/p\mathbb Z\)。
\end{proposition}
\begin{proof}
映射 \(\mathbb Z\rightarrow K\)、\(n\mapsto n1_K\) 的核为 \(m\mathbb Z\)，其中取 \(m\ge0\)。若 \(m=1\)，便有 \(1_K=0\)，与域非零矛盾，所以 \(m=0\) 或 \(m>1\)。若 \(m>1\) 是合数，可写成 \(m=ab\)，其中 \(1<a,b<m\)。此时 \(a,b\) 均不在核中，所以 \(a1_K\) 与 \(b1_K\) 都非零，但 \((a1_K)(b1_K)=m1_K=0\)，与域无零因子矛盾。因此 \(m=0\) 或 \(m=p\) 为素数。在后一情形，像同构于域 \(\mathbb F_p\)，每个子域必须包含它。在前一情形，像是 \(\mathbb Z\) 的副本，取其非零元素的逆元便得到 \(\mathbb Q\) 的副本；映射 \(a/b\mapsto(a1_K)(b1_K)^{-1}\) 良定义且单射。任一子域必须包含这些元素，所以所得子域最小，唯一性也随之成立。
\end{proof}
例如 \(\mathbb Q\subseteq\mathbb R\subseteq\mathbb C\) 是域扩张塔；\(\mathbb Z\subseteq\mathbb Q\) 则不是域扩张，因为 \(\mathbb Z\) 不是域。域的特征不会随扩张改变，因为单位及其整数倍均未改变。

\subsection{有限扩张与塔式公式}
\label{b1:subsec:1401:02}

设 \(L/K\) 为域扩张。以 \(L\) 的加法为加法，以域 \(L\) 中的乘法 \(K\times L\rightarrow L\)、\((c,\ell)\mapsto c\ell\) 为标量乘法，\(L\) 成为 \(K\) 向量空间\footnote{“向量空间”也称“线性空间”；参见席南华《基础代数》第二卷\cite[定义1.1, p. 1]{Xi2018BasicAlgebraII}。}。这里 \(c\in K\) 本来也是 \(L\) 的元素，所以标量乘法就是在 \(L\) 中作乘法；分配律、标量乘法结合律以及 \(1\ell=\ell\) 都继承自域 \(L\) 的运算。把这个向量空间的维数称为扩张 \(L/K\) 的\term[kuozhang-cishu]{扩张次数}[degree of an extension]，记为 \([L:K]\)。如果 \([L:K]\) 有限，则称 \(L/K\) 为\term[youxian-kuozhang]{有限扩张}[finite extension]；否则称为无限扩张。这里的有限指维数，不指底集合，例如 \(\mathbb Q(\sqrt2)/\mathbb Q\) 的次数为二，但两个域都是无限集。
\symidx{\([L:K]\)：域扩张的次数}

\ProseParagraphBreak
例如在域扩张塔 \(K\subseteq E\subseteq L\) 中，若 \(E/K\) 的一组基为 \(\{1,u\}\)，\(L/E\) 的一组基为 \(\{1,v,v^2\}\)，则任意 \(z\in L\) 唯一写成 \(e_0+e_1v+e_2v^2\)，其中 \(e_i\in E\)。再将每个 \(e_i\) 唯一写成 \(a_i+b_i u\)，其中 \(a_i,b_i\in K\)，便得到
\[
z=a_0+b_0u+a_1v+b_1uv+a_2v^2+b_2uv^2.
\]
两层坐标都唯一，所以 \(\{1,u,v,uv,v^2,uv^2\}\) 是 \(L/K\) 的一组基：三个 \(E\) 坐标各需两个 \(K\) 坐标，总共需 \(3\cdot2=6\) 个 \(K\) 坐标。下面的公式把这种逐层展开推广到任意两组基。

\ProseParagraphBreak
有限 Galois 理论的次数计算只涉及有限维时的整数公式。一般扩张也允许无限基，此时乘积基的大小用基数乘法表示。\secref{b1:sec:1403}比较无限超越基的大小时，还将使用无限基数的计数与比较。

\begin{theorem}[塔式公式]\label{b1:thm:tower-law}
设 \(K\subseteq E\subseteq L\) 为域扩张塔。若 \((u_i)_{i\in I}\) 为 \(E\) 在 \(K\) 上的一组向量空间基，\((v_j)_{j\in J}\) 为 \(L\) 在 \(E\) 上的一组向量空间基，则 \((u_iv_j)_{(i,j)\in I\times J}\) 为 \(L\) 在 \(K\) 上的一组向量空间基。因此
\[
[L:K]=[L:E]\,[E:K],
\]
无限维时右边理解为基数乘法。特别地，\(L/K\) 有限当且仅当 \(L/E\) 与 \(E/K\) 都有限。
\end{theorem}
\begin{proof}
任意 \(z\in L\) 是有限多个 \(v_j\) 的 \(E\) 线性组合；把这些系数分别写成有限多个 \(u_i\) 的 \(K\) 线性组合，即得 \(z\) 是有限多个 \(u_iv_j\) 的 \(K\) 线性组合。若有限和 \(\sum_{i,j}a_{ij}u_iv_j=0\)，其中 \(a_{ij}\in K\)，先按 \(j\) 收集。\(v_j\) 在 \(E\) 上线性无关，故每个 \(\sum_i a_{ij}u_i=0\)；再用 \(u_i\) 在 \(K\) 上线性无关，得所有 \(a_{ij}=0\)。这证明所列乘积是基。两个非零基数的乘积有限，恰好在二者均有限时发生。
\end{proof}
这一结果称为\term[tashi-gongshi]{塔式公式}[tower law]。扩张次数为一恰好等价于两域相等：若 \([L:K]=1\)，非零向量 \(1\) 本身就是 \(L/K\) 的一组基，因此 \(L=K\cdot1=K\)；反过来，\(K/K\) 以 \(1\) 为基，次数为一。有限扩张中，中间域的次数必整除总次数。因此，素数次数的扩张没有严格处于两端之间的中间域；这个判断不用知道扩张的具体生成元。

\subsection{单扩张及最小多项式}
\label{b1:subsec:1401:03}

\begin{definition}\label{b1:def:algebraic-element}
设 \(L/K\) 为域扩张，\(\alpha\in L\)。如果存在非零多项式 \(f\in K[x]\) 使 \(f(\alpha)=0\)，则称 \(\alpha\) 为 \(K\) 上的\term[daishu-yuan]{代数元}[algebraic element]，也说它在 \(K\) 上代数；否则称其为 \(K\) 上的\term[chaoyue-yuan]{超越元}[transcendental element]，也说它在 \(K\) 上超越。如果一个域扩张由单个元素 \(\alpha\) 生成，即形如 \(K(\alpha)/K\)，则称它为\term[dan-kuozhang]{单扩张}[simple extension]。
\end{definition}

求值同态
\[
\operatorname{ev}_{\alpha}:K[x]\longrightarrow L,\qquad
f\longmapsto f(\alpha)
\]
的像记为 \(K[\alpha]\)，即 \(\alpha\) 的全部 \(K\) 系数多项式值；它的核则收集了 \(\alpha\) 满足的全部多项式关系。代数情形下，这个核非零；由于 \(K[x]\) 是主理想整环，这个关系理想有唯一的首一生成元。
\symidx{\(\operatorname{ev}_{\alpha}\)：多项式的代入同态}
\symidx{\(K[\alpha]\)：元素 \(\alpha\) 的 \(K\) 系数多项式值组成的子环}

\begin{proposition}\label{b1:prop:minimal-polynomial-degree}
设 \(L/K\) 为域扩张，\(\alpha\in L\)。若 \(\alpha\) 在 \(K\) 上代数，则存在唯一首一不可约多项式 \(m_{\alpha,K}\in K[x]\)，使对每个 \(f\in K[x]\) 都有
\[
f(\alpha)=0\quad\Longleftrightarrow\quad m_{\alpha,K}\mid f.
\]
它称为 \(\alpha\) 在 \(K\) 上的\term[zuixiao-duoxiangshi]{最小多项式}[minimal polynomial]。若其次数为 \(d\)，则
\[
K(\alpha)=K[\alpha]\cong K[x]/(m_{\alpha,K}),\qquad
\bigl[K(\alpha):K\bigr]=d,
\]
且 \(1,\alpha,\ldots,\alpha^{d-1}\) 为一组基。若 \(\alpha\) 超越，则代入 \(x\mapsto\alpha\) 给出 \(K(\alpha)\cong K(x)\)，其中 \(K(x)\) 是多项式环 \(K[x]\) 的分式域。
\end{proposition}
\begin{proof}
代数情形下，求关系理想的首一生成元可以从次数最小的关系入手。在零化 \(\alpha\) 的非零多项式中取次数最小者，并除以首项系数，得到首一 \(m\)。对任意 \(f\) 作\cref{b1:prop:field-poly-division}中的带余除法 \(f=qm+r\)、\(\deg r<\deg m\)。若 \(f(\alpha)=0\)，则 \(r(\alpha)=0\)，最小性迫使 \(r=0\)。若 \(m=gh\) 而两因子均为正次数，则域中 \(g(\alpha)h(\alpha)=0\) 迫使一个因子零化 \(\alpha\)，再次矛盾。故 \(m\) 不可约；另一个满足条件的首一多项式与它相互整除，必相等。

代入同态 \(K[x]\rightarrow L\) 的像是 \(K[\alpha]\)，核为 \((m)\)，因此商环与 \(K[\alpha]\) 同构。要说明这个像已是域，目标是把每个非零 \(g(\alpha)\) 的逆也写成 \(\alpha\) 的多项式值。若 \(g(\alpha)\neq0\)，则 \(m\nmid g\)，由不可约性知 \(\gcd(m,g)=1\)。多项式 Euclid 算法给出 \(um+vg=1\)，代入后得 \(v(\alpha)g(\alpha)=1\)，恰好找到了这样的逆。所以 \(K[\alpha]\) 已是域，等于 \(K(\alpha)\)。带余除法给出所列幂的张成性；次数小于 \(d\) 的多项式不能零化 \(\alpha\)，给出线性无关性。

超越情形下代入同态的核为零，故 \(K[\alpha]\cong K[x]\)。取分式就把它延拓为 \(K(x)\rightarrow L\)，像恰是包含 \(K,\alpha\) 的最小子域，亦即 \(K(\alpha)\)。
\end{proof}
\symidx{\(m_{\alpha,K}\)：\(\alpha\) 在 \(K\) 上的最小多项式}
代数情形下，\(K[\alpha]=K(\alpha)\)，每个非零多项式值的逆仍是多项式值，取逆不增加元素。超越情形下，\(K[\alpha]\cong K[x]\) 只是多项式环，而 \(K(\alpha)\cong K(x)\) 允许多项式比；例如 \(1/\alpha\) 不能写成 \(\alpha\) 的多项式，否则清除分母就会得到非零多项式关系。

\ProseParagraphBreak
反过来，设 \(K\) 为域，\(q\in K[x]\) 为首一不可约多项式。要构造一个根，可以把 \(q(x)=0\) 作为新关系，取商环
\[
L=K[x]/(q),\qquad \alpha=x+(q).
\]
由于 \(K[x]\) 是主理想整环，\cref{b1:prop:pid-irreducible-prime,b1:prop:prime-max-quotient}说明 \(L\) 是域。常数映射 \(K\rightarrow L\) 是单射，因为正次数的 \(q\) 不可能整除非零常数。将 \(K\) 与其像识别后，在商环中直接计算
\[
q(\alpha)=q(x)+(q)=(q)=0_L.
\]
理想 \((q)\) 中的多项式被变为零，故 \(x\) 的剩余类确实成为 \(q\) 的根。又有 \(L=K[\alpha]=K(\alpha)\)；由\cref{b1:prop:minimal-polynomial-degree}，\(\alpha\) 在 \(K\) 上的最小多项式就是 \(q\)。这便从给定多项式构造了含有其根的域扩张。

\ProseParagraphBreak

例如在 \(\mathbb F_2[x]\) 中，二次多项式 \(q=x^2+x+1\) 在 \(0\) 与 \(1\) 处的值都为 \(1\)，故不可约。取 \(L=\mathbb F_2[x]/(q)\)、\(\alpha=x+(q)\)。按 \(q\) 带余除法，\(L\) 的四个元素为 \(0,1,\alpha,\alpha+1\)。这里特征为二，所以 \(-1=1\)、\(-\alpha=\alpha\)；从 \(q(\alpha)=0\) 得
\[
\alpha^2=-\alpha-1=\alpha+1,\qquad \alpha(\alpha+1)=1.
\]
所以 \(1,\alpha\) 是 \(L\) 在 \(\mathbb F_2\) 上的一组基，第二个等式也给出 \(\alpha^{-1}=\alpha+1\)。

\ProseParagraphBreak

若根已经位于某个扩张域中，同一结论则用于辨认该扩张。例如 \(x^2-2\) 在 \(\mathbb Q[x]\) 中不可约，因此 \(\mathbb Q(\sqrt2)\) 的元素唯一写成 \(a+b\sqrt2\)。其中非零元素的逆为
\[
(a+b\sqrt2)^{-1}=\frac{a-b\sqrt2}{a^2-2b^2}.
\]
分母若为零，且 \(b\neq0\)，就会得到有理数 \(a/b\) 的平方为二；若 \(b=0\)，则 \(a=0\)。所以对非零元素，所写分母确实非零。最小多项式依赖基域：\(\sqrt2\) 在 \(\mathbb Q(\sqrt2)\) 上的最小多项式已变为 \(x-\sqrt2\)。

\begin{proposition}\label{b1:prop:biquadratic-example}
域 \(L=\mathbb Q(\sqrt2,\sqrt3)\) 在 \(\mathbb Q\) 上的次数为四，基为 \(1,\sqrt2,\sqrt3,\sqrt6\)。元素 \(\theta=\sqrt2+\sqrt3\) 生成 \(L\)，其最小多项式为 \(x^4-10x^2+1\)。
\end{proposition}
\begin{proof}
若 \(\sqrt3=a+b\sqrt2\)、\(a,b\in\mathbb Q\)，平方后比较 \(1,\sqrt2\) 的系数得 \(ab=0\)、\(a^2+2b^2=3\)。\(b=0\) 会使 \(3\) 为有理平方，\(a=0\) 会使 \(3/2\) 为有理平方，均不可能，因为有理平方的素因子指数都为偶数。因此 \(x^2-3\) 在 \(\mathbb Q(\sqrt2)\) 上无根，二次塔的两层次数都是二。由\cref{b1:thm:tower-law}，总次数为四，所列四个乘积构成基。

令 \(\theta=\sqrt2+\sqrt3\)，先计算
\[
(\sqrt2+\sqrt3)(\sqrt3-\sqrt2)=3-2=1.
\]
所以 \(\theta^{-1}=\sqrt3-\sqrt2\)，并且
\[
\sqrt2=\frac{\theta-\theta^{-1}}2,\qquad
\sqrt3=\frac{\theta+\theta^{-1}}2.
\]
两个平方根均属于 \(\mathbb Q(\theta)\)，因而 \(L=\mathbb Q(\theta)\)。再由 \(\theta^2=5+2\sqrt6\) 得 \(P(\theta)=0\)，其中 \(P=x^4-10x^2+1\)。由 \([\mathbb Q(\theta):\mathbb Q]=4\) 及\cref{b1:prop:minimal-polynomial-degree}，最小多项式 \(m_{\theta,\mathbb Q}\) 的次数为四，而且它整除 \(P\)。两者次数相同且都首一，故 \(m_{\theta,\mathbb Q}=P\)。
\end{proof}

% !TeX root = ../../Book1.tex
% 纲要标题：### 14.2 代数扩张
% 纲要位置：计划纲要.md，第 508 行。

\section{代数扩张}
\label{b1:sec:1402}


% 纲要内容（仅作写作计划，尚非教材正文）：
%
% * 代数元构成子域
% * 有限扩张蕴含代数扩张
% * 代数扩张不必有限
%

两个分别满足多项式方程的元素相加、相乘后，是否仍满足多项式方程？直接消元往往使计算膨胀；有限维向量空间提供了更短的解释。本节先把有限生成的代数数据放进一个有限扩张，再由此处理一般代数扩张。

\subsection{代数元构成的子域}
\label{b1:subsec:1402:01}

逐个添入代数元时，下一步的基域已经包含前面添入的元素。扩大基域不会使原来的零化方程失效，最小多项式的次数因而只能下降；这使各步次数可以用最初基域上的次数统一估计。

\begin{lemma}\label{b1:lem:finite-algebraic-generators}
设 \(L/K\) 为域扩张，\(r\geq0\) 为整数，\(\alpha_1,\ldots,\alpha_r\in L\)。若这些元素都在 \(K\) 上代数，则 \(K(\alpha_1,\ldots,\alpha_r)/K\) 有限。更具体地，若 \(d_i=\deg m_{\alpha_i,K}\)，其次数至多为 \(d_1\cdots d_r\)，其中空乘积约定为一。
\end{lemma}
\begin{proof}
令 \(K_0=K\)、\(K_i=K(\alpha_1,\ldots,\alpha_i)\)。对每个 \(1\leq i\leq r\)，多项式 \(m_{\alpha_i,K}\) 也属于 \(K_{i-1}[x]\)，且仍然零化 \(\alpha_i\)。由\cref{b1:prop:minimal-polynomial-degree}，\(m_{\alpha_i,K_{i-1}}\) 整除这个多项式，整除关系在 \(K_{i-1}[x]\) 中理解。因此
\[
[K_i:K_{i-1}]
=\deg m_{\alpha_i,K_{i-1}}
\leq\deg m_{\alpha_i,K}=d_i.
\]
由\cref{b1:thm:tower-law}沿这个有限塔相乘，得所述估计。\(r=0\) 时扩张就是 \(K/K\)，次数为一。
\end{proof}

在域扩张 \(E/K\) 中取 \(a\in E\)。若 \(1,a,a^2,\ldots\) 的每个有限初段都线性无关，\(E\) 的维数就不可能有限。因此有限维会强迫这串幂满足某个有限的非平凡线性关系，而关系的系数正好给出零化 \(a\) 的多项式。

\begin{lemma}\label{b1:lem:finite-dimensional-algebraic}
设 \(E/K\) 为域扩张。若 \([E:K]=n<\infty\)，则每个 \(a\in E\) 都在 \(K\) 上代数，而且 \(\deg m_{a,K}\leq n\)。
\end{lemma}
\begin{proof}
在 \(n\) 维 \(K\) 向量空间 \(E\) 中，\(n+1\) 个向量 \(1,a,\ldots,a^n\) 线性相关。一个不全为零的相关系数列给出次数至多 \(n\) 的非零多项式 \(f\)，且 \(f(a)=0\)。最小多项式整除 \(f\)，故次数不超过 \(n\)。
\end{proof}

\begin{proposition}\label{b1:prop:algebraic-elements-subfield}
设 \(L/K\) 为域扩张。则 \(L\) 中在 \(K\) 上代数的全部元素组成一个包含 \(K\) 的子域。
\end{proposition}
\begin{proof}
每个 \(c\in K\) 被 \(x-c\) 零化。若 \(a,b\) 代数，\cref{b1:lem:finite-algebraic-generators}说明 \(E=K(a,b)\) 有限。\(a-b,ab\) 及 \(a\neq0\) 时的 \(a^{-1}\) 均在 \(E\) 内，由\cref{b1:lem:finite-dimensional-algebraic}也在 \(K\) 上代数。这正是包含 \(K\) 的子域所需的封闭性。
\end{proof}
这个证明给出方程存在性的结构解释：\(a-b,ab\) 以及非零 \(a\) 的逆元都落在同一个有限扩张里，所以各自满足某个非零的 \(K\) 系数多项式方程。证明不必从 \(a,b\) 的方程直接算出这些新方程。

\ProseParagraphBreak
例如所有实代数数构成 \(\mathbb R\) 的子域。对于 \(\theta=\sqrt2+\sqrt3\)，有限扩张 \(\mathbb Q(\sqrt2,\sqrt3)\) 先保证方程存在；若要找到具体方程，再用 \(\theta^2=5+2\sqrt6\) 平方消元，得到 \(\theta^4-10\theta^2+1=0\)。\cref{b1:prop:biquadratic-example}还说明这个四次关系就是最小多项式。

\subsection{代数扩张及其传递性}
\label{b1:subsec:1402:02}

\begin{definition}\label{b1:def:algebraic-extension}
设 \(L/K\) 为域扩张。如果 \(L\) 的每个元素都在 \(K\) 上代数，则称 \(L/K\) 为\term[daishu-kuozhang]{代数扩张}[algebraic extension]；否则称其为超越扩张。如果存在有限集 \(S\subseteq L\) 使 \(L=K(S)\)，则称 \(L/K\) 为\term[youxian-shengcheng-kuozhang]{有限生成扩张}[finitely generated extension]。
\end{definition}
由\cref{b1:lem:finite-dimensional-algebraic}，有限扩张必代数。它也必有限生成：若 \(v_1,\ldots,v_n\) 是 \(L\) 在 \(K\) 上的一组有限线性基，那么包含 \(K\) 及这些基向量的子域，必包含全部 \(K\) 线性组合；这些线性组合就是 \(L\) 的所有元素，所以 \(K(v_1,\ldots,v_n)=L\)。

\ProseParagraphBreak
反过来，若 \(L/K\) 既代数又有限生成，选取有限个域生成元后，每个生成元都代数，\cref{b1:lem:finite-algebraic-generators}就说明 \(L/K\) 有限。因此
\[
\text{有限扩张}
\quad\Longleftrightarrow\quad
\text{有限生成的代数扩张}.
\]
右端的两个条件必须同时成立；本节最后的例子将分别检验只保留其中一个条件时会发生什么。

\begin{theorem}\label{b1:thm:algebraic-tower}
设 \(K\subseteq E\subseteq L\) 为域扩张塔。则 \(L/K\) 代数当且仅当 \(E/K\) 与 \(L/E\) 都代数。
\end{theorem}
\begin{proof}
关键是反向蕴含：\(E/K\) 可能无限，不能直接把整个 \(E\) 当作有限的一层。但零化一个固定元素 \(a\) 的多项式只有有限多个系数，因此可以先把这些系数放进有限中间域 \(E_0\)，再使用 \(K\subseteq E_0\subseteq E_0(a)\) 这个两层有限塔。

若 \(L/K\) 代数，限制到 \(E\) 即知 \(E/K\) 代数；把零化元素的 \(K\) 系数多项式视为 \(E\) 系数多项式，即知 \(L/E\) 代数。

反过来，设 \(E/K\) 与 \(L/E\) 都代数。给定 \(a\in L\)，取零化它的非零多项式 \(f=c_0+c_1x+\cdots+c_dx^d\in E[x]\)，并令 \(E_0=K(c_0,\ldots,c_d)\)。因每个 \(c_i\) 在 \(K\) 上代数，\cref{b1:lem:finite-algebraic-generators}说明 \(E_0/K\) 有限。原来的 \(f\) 已经属于 \(E_0[x]\)，且仍是非零多项式，所以 \(a\) 在 \(E_0\) 上代数，由\cref{b1:prop:minimal-polynomial-degree}知 \(E_0(a)/E_0\) 也有限。由\cref{b1:thm:tower-law}，\(E_0(a)/K\) 有限，再由\cref{b1:lem:finite-dimensional-algebraic}知 \(a\) 在 \(K\) 上代数。任意 \(a\) 均如此，故 \(L/K\) 代数。
\end{proof}
这里的 \(E_0\) 可以随 \(a\) 及所选方程而改变，不要求所有元素所需的系数同时落在一个固定有限子域中。这使论证也适用于 \(E/K\) 无限的情形。

\subsection{无限代数扩张}
\label{b1:subsec:1402:03}

\begin{proposition}\label{b1:prop:infinite-algebraic-root-tower}
对整数 \(n\geq1\)，在 \(\mathbb R\) 中取正实根 \(a_n=2^{1/2^n}\)，令
\[
K_n=\mathbb Q(a_n),\qquad L=\bigcup_{n\geq1}K_n.
\]
则 \(L/\mathbb Q\) 是无限代数扩张，而且不是有限生成扩张。
\end{proposition}
\begin{proof}
因为 \(a_n=a_{n+1}^2\)，这些域递增。任意两个元素落在同一个 \(K_n\) 中，所以其差、积及非零元素的逆也落在该域中；因此 \(L\) 是域。每个 \(a_n\) 满足 \(x^{2^n}-2=0\)，故 \(K_n/\mathbb Q\) 有限；每个 \(L\) 中的元素属于一个这样的有限扩张，因此在 \(\mathbb Q\) 上代数。

另一方面，\(x^{2^n}-2\) 对素数二满足\cref{b1:thm:eisenstein}的条件，故 \([K_n:\mathbb Q]=2^n\)。如果 \([L:\mathbb Q]\) 有限，\cref{b1:thm:tower-law}就会使它至少为每个 \(2^n\)，不可能。

若 \(L/\mathbb Q\) 有限生成，每个生成元又都在 \(\mathbb Q\) 上代数，\cref{b1:lem:finite-algebraic-generators}便会使 \(L/\mathbb Q\) 有限，与刚才的结论矛盾。因此它也不是有限生成扩张。
\end{proof}
这里每个 \(K_n/\mathbb Q\) 都有限，每个 \(a\in L\) 也都落入某个有限层；但层数可以随 \(a\) 改变，全体元素不能同时落入一个固定有限中间域。

\ProseParagraphBreak
单独的有限生成或代数条件都不能推出有限。若 \(t\) 是 \(K\) 上的不定元，\(K(t)/K\) 只需 \(t\) 一个元素生成，却不是代数扩张，也不是有限扩张，因为 \(1,t,t^2,\ldots\) 线性无关。上述根塔 \(L/\mathbb Q\) 则代数，却既非有限也非有限生成。这两个反例分别说明，不能从“有限生成的代数扩张”中去掉代数或有限生成的条件。

\ProseParagraphBreak
有限中间域未必排成一条升链，但它们可以组成有向族。这里“有向”的含义是：任取有限多个有限中间域，都能找到一个仍然有限的中间域共同包含它们。这样，原来分别位于不同有限层的元素就能移到同一层中计算。

\begin{proposition}\label{b1:prop:algebraic-directed-union}
设 \(L/K\) 为代数域扩张。则 \(L\) 是其有限中间扩张的并；任意有限多个这样的中间域都包含于另一个有限中间扩张。
\end{proposition}
\begin{proof}
每个 \(a\in L\) 属于有限扩张 \(K(a)\)。若 \(E_1,\ldots,E_s\) 为有限中间域，分别选取它们在 \(K\) 上的有限基 \(B_1,\ldots,B_s\)，并令 \(E=K(B_1\cup\cdots\cup B_s)\)。合在一起的基向量是有限多个代数元，故由\cref{b1:lem:finite-algebraic-generators}，\(E/K\) 仍有限。对每个 \(i\)，\(E\) 包含 \(K\) 和 \(B_i\)，所以包含 \(B_i\) 的全部 \(K\) 线性组合，也就是整个 \(E_i\)。而所用生成元都在 \(L\) 中，故 \(E\subseteq L\)，这给出共同的有限中间域。
\end{proof}
例如计算分别来自有限中间域 \(E_1,E_2\) 的两个元素的差或积时，先选共同包含这两个域的有限中间域 \(E\)，运算就能在 \(E\) 内完成；非零元素的逆也留在该域中。因此每个只涉及有限多个元素的计算都可移到某个有限层。处理整个扩张的自同构或子域结构时，则必须另外检查各有限层之间的相容性；有限层的结论不能不经说明便推到无限层。

% !TeX root = ../../Book1.tex
% 纲要标题：### 14.3 超越扩张
% 纲要位置：计划纲要.md，第 514 行。

\section{超越扩张}
\label{b1:sec:1403}


% 纲要内容（仅作写作计划，尚非教材正文）：
%
% * 代数独立与超越基
% * 交换引理及超越次数的不变性
% * 有限生成扩张与纯超越扩张
% * 超越次数的塔式公式
%

一个超越元没有单变量多项式关系，几个超越元之间却仍可能有关系。例如 \(t\) 与 \(t^2\) 分别在 \(K\) 上超越，但满足 \(Y-X^2=0\)。要衡量扩张中真正独立的参数数目，必须同时考察有限多个元素的多项式关系。

\ProseParagraphBreak
有限生成扩张可从已给生成元中逐个选取独立参数。一般超越基的存在性则使用 Zorn 引理；比较无限超越基的大小，还需对有限支集作基数计数。

\subsection{代数独立与超越基}
\label{b1:subsec:1403:01}

\begin{definition}\label{b1:def:transcendence-basis}
设 \(L/K\) 为域扩张，\(T\subseteq L\)。如果对任意整数 \(r\geq1\)、任意互异的 \(t_1,\ldots,t_r\in T\) 以及任意 \(f\in K[X_1,\ldots,X_r]\)，等式 \(f(t_1,\ldots,t_r)=0\) 都蕴含 \(f=0\)，则称 \(T\) 在 \(K\) 上\term[daishu-duli]{代数独立}[algebraic independence]。如果 \(T\) 在 \(K\) 上代数独立，并且 \(L/K(T)\) 是代数扩张，则称 \(T\) 为 \(L/K\) 的\term[chaoyue-ji]{超越基}[transcendence basis]。空集视为代数独立，并约定 \(K(\varnothing)=K\)。
\end{definition}

\begin{lemma}\label{b1:lem:independence-adjoining}
设 \(L/K\) 为域扩张，\(T\subseteq L\) 在 \(K\) 上代数独立，且 \(a\in L\setminus T\)。则 \(T\cup\{a\}\) 在 \(K\) 上代数独立，当且仅当 \(a\) 在 \(K(T)\) 上超越。因此，\(L/K\) 的超越基恰是 \(L\) 中按包含关系极大的 \(K\) 上代数独立子集。
\end{lemma}
\begin{proof}
\(K(T)\) 是由 \(T\) 中有限多个元素的多项式比组成的域。若 \(a\) 满足 \(K(T)\) 系数的多项式方程，清除有限多个分母，就得到 \(a\) 与 \(T\) 中有限多个元素之间的非零多项式关系；多项式仍非零，是因为 \(T\) 代数独立。反过来，把一个这样的关系按 \(a\) 的幂收集，至少一个系数是 \(T\) 中元素的非零多项式值，因而非零；这便给出 \(a\) 在 \(K(T)\) 上的代数方程。极大性的等价刻画随即得到。
\end{proof}

对于有限生成扩张，可以直接选择超越基。若 \(L=K(s_1,\ldots,s_n)\)，依次检查这些生成元；只在 \(s_i\) 对已选元素生成的域仍超越时，才把它选入集合 \(T\)。由\cref{b1:lem:independence-adjoining}，\(T\) 始终代数独立。每个未选入的 \(s_i\) 在当时的参数域上代数，因而在最终的 \(K(T)\) 上仍代数；由\cref{b1:lem:finite-algebraic-generators}，它们有限个生成的扩张 \(L/K(T)\) 也代数，所以 \(T\) 是 \(L/K\) 的超越基。一般扩张未必有有限生成集可供逐项检查，下面的存在定理因而采用 Zorn 引理。

\begin{theorem}\label{b1:thm:transcendence-basis}
每个域扩张 \(L/K\) 都有超越基。更一般地，任意在 \(K\) 上代数独立的子集 \(A\subseteq L\) 都可扩充成 \(L/K\) 的超越基。
\end{theorem}
\begin{proof}
考虑包含 \(A\) 的代数独立子集，按包含排序。它包含 \(A\)，因而非空。任一非空链的并仍代数独立：一个多项式关系只涉及有限多个元素，它们已经同属链中的某一个成员。这个并仍包含 \(A\)，故是该链的上界；空链则以 \(A\) 为上界。于是 Zorn 引理（\cref{b1:thm:A-zorn}）给出极大元 \(T\)。由\cref{b1:lem:independence-adjoining}，每个不在 \(T\) 中的元素都在 \(K(T)\) 上代数；\(T\) 中元素本来就在该域中，所以 \(L/K(T)\) 代数。
\end{proof}

\subsection{交换引理及超越次数的不变性}
\label{b1:subsec:1403:02}

\begin{lemma}[交换引理]\label{b1:lem:transcendence-exchange}
设 \(L/F\) 为域扩张，\(x,y\in L\)。若 \(x\) 在 \(F(y)\) 上代数而在 \(F\) 上超越，则 \(y\) 在 \(F(x)\) 上代数。
\end{lemma}
\begin{proof}
把 \(x\) 的一个 \(F(y)\) 系数方程清除分母，得到 \(P(x,y)=0\)，其中 \(P(X,Y)\in F[X,Y]\) 非零。把它按 \(Y\) 展开为 \(\sum_j q_j(X)Y^j\)。由于 \(x\) 在 \(F\) 上超越，每个非零 \(q_j\) 都满足 \(q_j(x)\neq0\)，所以 \(P(x,Y)\) 是 \(F(x)[Y]\) 的非零多项式。它以 \(y\) 为根，故结论成立。
\end{proof}
该结论称为\term[jiaohuan-yinli]{交换引理}[exchange lemma]。它允许用一个新的独立参数替换旧参数，但替换时必须保证新参数在其余已保留参数上超越。

\begin{lemma}\label{b1:lem:finite-transcendence-bound}
设 \(L/K\) 为域扩张，\(b_1,\ldots,b_m\in L\)。若 \(L/K(b_1,\ldots,b_m)\) 代数，则 \(L\) 中任意在 \(K\) 上代数独立的子集至多含 \(m\) 个元素。这里不要求 \(b_1,\ldots,b_m\) 本身代数独立。
\end{lemma}
\begin{proof}
反设 \(a_1,\ldots,a_{m+1}\in L\) 在 \(K\) 上代数独立。设已选取其中的 \(a_1,\ldots,a_r\)，令 \(F=K(a_1,\ldots,a_r)\)，并以 \(D\) 表示尚未替换的 \(b_i\) 构成的集合。归纳假设是 \(\lvert D\rvert\leq m-r\) 且 \(L/E\) 代数，其中 \(E=F(D)\)；\(r=0\) 时这正是题设。令 \(a=a_{r+1}\)。它在 \(F\) 上超越，却在 \(E\) 上代数，因此 \(D\) 非空。在 \(D\) 中取一个包含极小的非空子集 \(C\)，使 \(a\) 在 \(F(C)\) 上代数。任取 \(b\in C\)，若删去 \(b\) 后 \(a\) 仍代数，就与 \(C\) 的极小性矛盾。因此 \(a\) 在 \(F\bigl(C\setminus\{b\}\bigr)\) 上仍超越；这正是交换引理所需的假设。

对底域 \(F\bigl(C\setminus\{b\}\bigr)\) 应用\cref{b1:lem:transcendence-exchange}，得 \(b\) 在 \(F\bigl(C\setminus\{b\},a\bigr)\) 上代数。令
\[
E'=F\bigl(D\setminus\{b\},a\bigr),\qquad
M=F(D,a)=E(a)=E'(b).
\]
由于 \(a\) 在 \(E\) 上代数，\(M/E\) 代数。交换所得的方程在更大的系数域 \(E'\) 上仍有效，故 \(b\) 在 \(E'\) 上代数，即 \(M/E'\) 代数。又因 \(E\subseteq M\subseteq L\)，归纳假设给出 \(L/M\) 代数。旧参数域 \(E\) 与新参数域 \(E'\) 未必互相包含，不能把二者直接排成一条域塔；共同的上层域 \(M\) 给出了如下包含关系：
\[
\begin{tikzcd}[row sep=large, column sep=large]
& L & \\
& M \arrow[u, hook, "\text{代数}"] & \\
E \arrow[ur, hook, "\text{代数}"] && E' \arrow[ul, hook, "\text{代数}"']
\end{tikzcd}
\]
沿右侧的 \(E'\subseteq M\subseteq L\) 应用\cref{b1:thm:algebraic-tower}，得到 \(L/E'\) 代数。这样便用 \(a\) 替换了 \(b\)，同时保持了归纳所需的代数性。每次替换使 \(D\) 少一个元素，故至多 \(m\) 次后它为空。此时 \(L/F\) 代数，但 \(a_{r+1}\) 在 \(F\) 上超越，矛盾。
\end{proof}

\begin{theorem}\label{b1:thm:transcendence-cardinality}
设 \(L/K\) 为域扩张。任意两个 \(L/K\) 的超越基都具有相同基数；这个共同基数称为扩张 \(L/K\) 的\term[chaoyue-cishu]{超越次数}[transcendence degree]，记为 \(\operatorname{trdeg}_K L\)。
\end{theorem}
\begin{proof}
设 \(A,B\) 为超越基。若其中一个有限，不妨设 \(B\) 有限。\cref{b1:lem:finite-transcendence-bound}给出 \(\lvert A\rvert\leq\lvert B\rvert\)，于是 \(A\) 也有限；交换二者再用该引理即得等号。

现在设二者均无限。每个 \(a\in A\) 在 \(K(B)\) 上代数，其一个代数方程只使用 \(B\) 中有限多个元素的有理函数，故可选有限子集 \(B_a\subseteq B\)，使 \(a\) 在 \(K(B_a)\) 上代数。对任意有限集 \(C\subseteq B\)，令
\[
A_C=\{\,a\in A\mid B_a=C\,\},\qquad M_C=K(C)(A_C).
\]
\(A_C\) 的每个元素都在 \(K(C)\) 上代数；由\cref{b1:prop:algebraic-elements-subfield}，\(M_C/K(C)\) 也代数。把\cref{b1:lem:finite-transcendence-bound}用于 \(M_C/K\)，以 \(C\) 为至多 \(\lvert C\rvert\) 个参数，便得到 \(\lvert A_C\rvert\leq\lvert C\rvert\)，因为 \(A_C\subseteq A\) 仍在 \(K\) 上代数独立。

由\cref{b1:prop:finite-support-cardinality}，无限集 \(B\) 的有限子集总数为 \(\lvert B\rvert\)，而每个 \(A_C\) 都有限。因此对应 \(a\mapsto B_a\) 给出 \(\lvert A\rvert\leq\lvert B\rvert\)。交换 \(A,B\) 得反向不等式，再由\cref{b1:prop:cantor-bernstein}得到等号。为各 \(a\) 同时选取有限支集 \(B_a\)，以及这里的无限基数计数，均采用\cref{b1:def:A-choice}中的选择公理；它与选取超越基时使用的选择原则一致。
\end{proof}
\symidx{\(\operatorname{trdeg}_K L\)：扩张的超越次数}
最后一段不能仅用有限次交换代替：无限基的大小是基数问题，有限个参数的局部估计必须结合有限支集的计数，才能得到所需结论。

\subsection{有限生成扩张与纯超越扩张}
\label{b1:subsec:1403:03}

\begin{definition}\label{b1:def:purely-transcendental}
设 \(L/K\) 为域扩张。如果存在 \(K\) 上代数独立的子集 \(T\subseteq L\)，使 \(L=K(T)\)，则称 \(L/K\) 为\term[chun-chaoyue-kuozhang]{纯超越扩张}[purely transcendental extension]。
\end{definition}
对有限集 \(T=\{t_1,\ldots,t_r\}\)，代入给出 \(K(T)\cong K(X_1,\ldots,X_r)\)。一般超越扩张未必纯超越；超越基只保证以参数域 \(K(T)\) 为新基域后，\(L/K(T)\) 是代数扩张。

\begin{proposition}\label{b1:prop:finitely-generated-field-structure}
设 \(L/K\) 为有限生成域扩张。则存在有限超越基 \(T\subseteq L\)，使 \(L/K(T)\) 有限。因而每个有限生成扩张都可以分成一个有限参数的纯超越扩张和其后的有限代数扩张。
\end{proposition}
\begin{proof}
从有限生成集 \(S\) 中按上文的方法逐个选取，得到超越基 \(T\subseteq S\)。每个 \(s\in S\setminus T\) 在 \(K(T)\) 上代数；\(S\setminus T\) 有限，故\cref{b1:lem:finite-algebraic-generators}说明 \(L=K(T)\bigl(S\setminus T\bigr)\) 在 \(K(T)\) 上有限。
\end{proof}

\begin{proposition}[超越次数的塔式公式]\label{b1:prop:transcendence-tower}
设 \(K\subseteq E\subseteq L\) 为域扩张塔。则下式中的加法为基数加法：
\begin{equation}\label{b1:eq:transcendence-tower}
\operatorname{trdeg}_K L=\operatorname{trdeg}_K E+\operatorname{trdeg}_E L.
\end{equation}
\end{proposition}
\begin{proof}
取 \(E/K\) 的超越基 \(T\) 及 \(L/E\) 的超越基 \(U\)。\(T\subseteq E\)，而 \(U\) 的任何元素都不属于 \(E\)，否则它在 \(E\) 上代数；故 \(T\cap U=\varnothing\)。

先证 \(T\cup U\) 在 \(K\) 上代数独立。只涉及 \(T\) 或只涉及 \(U\) 的关系已由各自的独立性排除。一般关系涉及有限多个 \(t_1,\ldots,t_r\in T\) 及 \(u_1,\ldots,u_s\in U\)，按 \(u_j\) 的单项式收集后可写成
\[
\sum_{\beta}P_\beta(t_1,\ldots,t_r)
u_1^{\beta_1}\cdots u_s^{\beta_s}=0,
\qquad P_\beta\in K[X_1,\ldots,X_r],
\]
其中 \(\beta\in\mathbb Z_{\geq0}^s\)，只有有限多个 \(P_\beta\) 非零。各系数 \(P_\beta(t_1,\ldots,t_r)\) 属于 \(E\)；由 \(U\) 在 \(E\) 上代数独立，它们全为零。再由 \(T\) 在 \(K\) 上代数独立，每个 \(P_\beta\) 都是零多项式，故原关系只能是零关系。

又因 \(T\) 是 \(E/K\) 的超越基，\(E/K(T)\) 代数。\(E\) 中每个元素在更大的 \(K(T,U)\) 上仍代数，所以 \(E(U)=K(T,U)(E)\) 在 \(K(T,U)\) 上代数。\(L/E(U)\) 也因 \(U\) 是 \(L/E\) 的超越基而代数；由\cref{b1:thm:algebraic-tower}，\(L/K(T,U)\) 代数。因此 \(T\cup U\) 是 \(L/K\) 的超越基。由于 \(T\cap U=\varnothing\)，其基数为 \(\lvert T\rvert+\lvert U\rvert\)，即得所述等式。
\end{proof}

有限代数扩张的次数沿塔相乘；超越次数沿塔按基数相加。

% !TeX root = ../../Book1.tex
% 纲要标题：### 14.4 函数域与超越次数
% 纲要位置：计划纲要.md，第 520 行。

\section{函数域与超越次数}
\label{b1:sec:1404}


% 纲要内容（仅作写作计划，尚非教材正文）：
%
% * 有理函数域与代数函数域
% * 次数与超越次数的不同角色；剩余有限次数随超越基的选择而变
% * 第三卷维数理论的域论背景
%
% ---
%

从多项式环取分式域，就得到可以进行除法的函数表达式。变量之间若有代数关系，取分式以后关系仍然存在。本节用几个可直接计算的例子说明：扩张次数描述代数方程的复杂度，超越次数则记录独立参数的数目。

\subsection{有理函数域与代数函数域}
\label{b1:subsec:1404:01}

\begin{definition}\label{b1:def:function-field}
设 \(L/k\) 为域扩张。如果 \(L/k\) 有限生成，则称 \(L\) 为 \(k\) 上的\term[hanshu-yu]{函数域}[function field]。本书允许函数域的超越次数为零，因而有限代数扩张也属于函数域。设 \(r\geq0\) 为整数；如果 \(t_1,\ldots,t_r\in L\) 在 \(k\) 上代数独立，则称 \(k(t_1,\ldots,t_r)\) 为 \(k\) 上的\term[youli-hanshu-yu]{有理函数域}[rational function field]；有理函数域的有限代数扩张通常称为\term[daishu-hanshu-yu]{代数函数域}[algebraic function field]。若 \(L\) 中所有在 \(k\) 上代数的元素都属于 \(k\)，则称 \(k\) 在 \(L\) 中相对代数闭。本定义不要求这一条件；它也不同于要求 \(k\) 自身是代数闭域。
\end{definition}
由\cref{b1:prop:finitely-generated-field-structure}，每个函数域都是某个有理函数域的有限扩张。例如 \(\mathbb Q(\sqrt2)/\mathbb Q\) 的超越次数为零，仍符合本书的函数域定义；此时超越基为空，参数域就是 \(\mathbb Q\)。

\ProseParagraphBreak
相对代数闭只比较 \(k\) 与当前的扩域 \(L\)。例如设 \(t\) 在 \(\mathbb Q\) 上超越，则 \(\mathbb Q\) 在 \(\mathbb Q(t)\) 中相对代数闭：其中在 \(\mathbb Q\) 上代数的有理函数只有有理常数，\cref{b1:ex:14-10}给出了这一事实的核验。但 \(\mathbb Q\) 本身不是代数闭域，因为 \(x^2-2\) 没有有理根。前者只问所给扩域是否带来新的代数元，后者要求 \(k[x]\) 中每个非常数多项式都有 \(k\) 中的根。

\ProseParagraphBreak
若整环 \(A\) 包含 \(k\)，且 \(A=k[a_1,\ldots,a_n]\)，其分式域是函数域。反过来，若 \(L=k(a_1,\ldots,a_n)\)，先取 \(L\) 中的子环 \(A=k[a_1,\ldots,a_n]\)。它位于域 \(L\) 中，因而是整环；由\cref{b1:prop:fraction-in-ambient-field}，\(\operatorname{Frac}(A)\) 正是 \(L\) 中包含 \(A\) 的最小子域，即
\[
L=k(a_1,\ldots,a_n)=\operatorname{Frac}\bigl(k[a_1,\ldots,a_n]\bigr).
\]
写 \(k[a_1,\ldots,a_n]\) 时，元素由生成元的 \(k\) 系数多项式给出；写 \(k(a_1,\ldots,a_n)\) 时，则还允许非零分母。二者未必相等：若 \(t\) 超越，\(1/t\in k(t)\)，却不在 \(k[t]\) 中，因为多项式 \(tq(t)\) 的常数项为零，不可能等于 \(1\)。

\ProseParagraphBreak
例如，设 \(k\) 为域、\(t\) 在 \(k\) 上超越，则 \(k[t^2,t^3]\) 的分式域仍为 \(k(t)\)，因为 \(t=t^3/t^2\)。所以不同的环可能有同一个函数域。还可考虑
\[
L=k(t)(u),\qquad u^2=t^3-t.
\]
多项式 \(x^2-(t^3-t)\) 在 \(k(t)[x]\) 中不可约。为证明这一点，假设 \(t^3-t=(A/B)^2\)，其中 \(A,B\in k[t]\) 均非零，并已将分式约成 \(\gcd(A,B)=1\) 的形式。清除分母后有
\[
A^2=(t^3-t)B^2=t(t-1)(t+1)B^2.
\]
在 \(k[t]\) 的唯一分解中，\(t\) 在 \(t^3-t\) 中恰出现一次。因此右边的因子 \(t\) 出现次数为一加上它在 \(B\) 中出现次数的两倍，是奇数；左边的出现次数则是它在 \(A\) 中出现次数的两倍，是偶数，矛盾。所以 \(t^3-t\) 不是 \(k(t)\) 中的平方。二次多项式无根便不可约，故可用不可约多项式的商域构造 \(L\)，并有 \(\bigl[L:k(t)\bigr]=2\)。

\ProseParagraphBreak
这给出一个带有两个生成元 \(t,u\)、一个关系的函数域，但不能将超越次数机械地算作“生成元数减方程数”。方程可以重复或由其他方程推出，生成元也可以冗余；独立参数的数目须由代数独立性判断。这里 \(t\) 超越，而 \(u\) 在 \(k(t)\) 上代数，才是超越次数为一的理由。刚才算出的次数二是相对于参数域 \(k(t)\) 的次数，并未判断是否存在另一参数 \(s\) 使 \(L=k(s)\)。

\subsection{次数与超越次数}
\label{b1:subsec:1404:02}

设 \(t\) 在 \(k\) 上超越，\(n\geq1\)。则 \(t^n\) 仍超越，否则将零化 \(t^n\) 的非零多项式中的变量换成 \(x^n\)，会得到零化 \(t\) 的非零多项式。并且
\begin{equation}\label{b1:eq:rational-power-degree}
\bigl[k(t):k(t^n)\bigr]=n.
\end{equation}
事实上，\(t\) 被 \(x^n-t^n\in k(t^n)[x]\) 零化，所以次数至多为 \(n\)。若 \(1,t,\ldots,t^{n-1}\) 存在非平凡的 \(k(t^n)\) 线性关系，清除分母得到
\[
\sum_{i=0}^{n-1}p_i(t^n)t^i=0,\qquad p_i\in k[x].
\]
清除分母后至少有一个 \(p_i\ne0\)。写 \(p_i(X)=\sum_jc_{ij}X^j\)，则上式变为
\[
\sum_{i=0}^{n-1}\sum_j c_{ij}t^{nj+i}=0.
\]
第 \(i\) 组中的幂指数模 \(n\) 余 \(i\)；在 \(0\le i<n\) 的范围内，不同的 \((i,j)\) 给出不同的指数 \(nj+i\)，所以不同组没有同次幂，不能相互抵消。\(t\) 超越使每个系数 \(c_{ij}=0\)，进而所有 \(p_i=0\)，与关系非平凡矛盾，证明了线性无关性。次数因而恰为 \(n\)，而 \(x^n-t^n\) 是 \(n\) 次首一零化多项式；由\cref{b1:prop:minimal-polynomial-degree}，它正是 \(t\) 在 \(k(t^n)\) 上的最小多项式，在 \(k(t^n)[x]\) 中不可约。

\ProseParagraphBreak
固定同一个纯超越扩张 \(L=k(t)/k\)。由上面的证明，\(\{t\}\)、\(\{t^2\}\) 和 \(\{t^n\}\) 都是它的超越基：这些元素在 \(k\) 上超越，且 \(t\) 在各自生成的参数域上代数。采用超越基 \(\{t\}\) 时，参数域为 \(k(t)=L\)，剩余次数 \([L:k(t)]=1\)；采用 \(\{t^2\}\) 时，参数域为 \(k(t^2)\)，剩余次数 \([L:k(t^2)]=2\)；一般地，对 \(n\ge1\)，采用 \(\{t^n\}\) 时，参数域为 \(k(t^n)\)，剩余次数 \([L:k(t^n)]=n\)。

\ProseParagraphBreak
扩张 \(L/k\) 的超越次数始终为一，剩余有限次数却随参数域而变。特别地，\([L:k(t^2)]=2\) 不能用来断定 \(L/k\) 不是纯超越扩张；\(L=k(t)\) 已经给出了反例。

\ProseParagraphBreak
当 \(\operatorname{char}k=p\) 且 \(n=p\) 时，次数仍为 \(p\)，\(x^p-t^p\) 在 \(k(t^p)[x]\) 中仍是不可约的最小多项式。把系数域扩大到 \(k(t)\) 后，才有分解 \(x^p-t^p=(x-t)^p\)，其中 \(t\) 已经属于新的系数域。不可约性必须注明系数域；这个分解也说明它只有一个不同的根，根的个数与扩张次数须加以区分，\chapref{b1:ch:16}将解释其差别。

\ProseParagraphBreak
上一小节的 \(L=k(t)(u)\) 同样满足 \(\operatorname{trdeg}_k L=1\)，因为 \(L/k(t)\) 代数，而 \(\{t\}\) 仍独立，故 \(\{t\}\) 是 \(L/k\) 的超越基。关系 \(u^2=t^3-t\) 表明 \(u\) 在这个参数域上代数，所以两个生成元只有一个独立参数。

\subsection{维数理论的域论背景}
\label{b1:subsec:1404:03}

把 \(t_1,\ldots,t_r\) 当作自由参数时，多项式环 \(k[t_1,\ldots,t_r]\) 没有参数之间的多项式关系，其分式域的超越次数就是 \(r\)。现在取 \(k[X_1,\ldots,X_n]\) 的素理想 \(\mathfrak p\)，令
\[
A=k[X_1,\ldots,X_n]/\mathfrak p.
\]
由\cref{b1:prop:prime-max-quotient}，素理想条件保证 \(A\) 是整环，因而可以取分式域。非零常数是多项式环的单位，不能属于真理想 \(\mathfrak p\)，所以 \(k\) 自然嵌入 \(A\)。剩余类 \(x_i\) 生成函数域 \(\operatorname{Frac}(A)=k(x_1,\ldots,x_n)\)。从这组生成元中抽出的超越基刻画其中相互独立的参数；其余生成元在这些参数上代数。

\ProseParagraphBreak
例如 \(k[X,Y]/(Y-X^2)\cong k[t]\)，同构由 \(X\mapsto t,Y\mapsto t^2\) 给出：用 \(Y=X^2\) 可把每个剩余类化为 \(X\) 的多项式，而这样的多项式映到零只能本来为零。商环中 \(Y\) 的剩余类已由 \(X\) 的剩余类平方得到，因此 \(Y\) 是这组生成元中的冗余项，可以消去。所得函数域 \(k(t)\) 只有一个独立参数。一般的方程也可能冗余，商环可能不再是整环，所以增加一个方程后，独立参数的变化要根据具体理想判断。

\ProseParagraphBreak
第三卷第5章将建立整扩张与 Noether 正规化，第9章将定义 Krull 维数\footnote{克鲁尔（Wolfgang Krull，1899-1971），德国数学家，发展交换环与理想理论，定义了以其名字命名的维数。}，并在域上有限生成整环的情形证明维数等于函数域的超越次数。这里尚未定义那种环的维数，因而只把超越次数视为已建立的域论不变量。


\begin{chaptersummary}
最小多项式把单个代数元的全部多项式关系压缩成一个首一不可约多项式，其次数等于相应单扩张的维数。域扩张塔中的基可以逐层相乘，得到次数的乘法公式。有限个代数元总落在有限扩张中，因此代数元对域运算封闭，代数性也沿扩张塔传递；无限代数扩张则可看作有限中间扩张的有向并。

\ProseParagraphBreak
超越基保留相互独立的参数，使剩余扩张变为代数扩张。交换引理与有限支集计数保证其基数不随选择改变。有限生成扩张因此分成纯超越部分和有限代数部分；超越次数由原扩张决定，剩余有限次数却可能随参数的选择而改变。

\ProseParagraphBreak
本章主要研究加入给定元素后的次数与代数性。下一章转向同时容纳一个多项式全部根的扩张，讨论这种扩张的构造及其同构意义下的唯一性。
\end{chaptersummary}

\BookExercises

\begin{exercise}\label{b1:ex:14-1}
证明域 \(K\) 的每个自同构都逐点固定其素域。若 \(K\) 有限，证明其元素个数必为 \(p^n\)，其中 \(p=\operatorname{char}K\)、\(n\geq1\)。
\end{exercise}
\begin{solution}
自同构固定 \(1\)，因而固定所有整数倍 \(m1\)。特征 \(p\) 时这些元素已组成素域；特征零时还保持逆元，所以固定所有 \((a1)(b1)^{-1}\)，即固定 \(\mathbb Q\) 的副本。有限域不可能含无限素域 \(\mathbb Q\)，故特征为素数 \(p\)。把 \(K\) 看成 \(\mathbb F_p\) 向量空间，因其底集有限可取有限基，设长度为 \(n\)。每个元素与唯一的 \(n\) 元坐标组对应，每个坐标有 \(p\) 种可能，故元素总数为 \(p^n\)。
\end{solution}

\begin{exercise}\label{b1:ex:14-2}
设 \(a^3=2\)。在 \(\mathbb Q(a)\) 中把 \((1+a)^{-1}\) 写成 \(1,a,a^2\) 的有理线性组合，并求 \(1+a\) 的最小多项式。
\end{exercise}
\begin{solution}
\((1+a)(1-a+a^2)=1+a^3=3\)，故逆元为 \((1-a+a^2)/3\)。若 \(b=1+a\)，则 \((b-1)^3-2=0\)，即 \(b^3-3b^2+3b-3=0\)。\(x^3-2\) 由 Eisenstein 判据不可约，平移 \(x\mapsto x-1\) 是 \(\mathbb Q[x]\) 的自同构，保持不可约性。因此 \(x^3-3x^2+3x-3\) 就是 \(b\) 的最小多项式。
\end{solution}

\begin{exercise}\label{b1:ex:14-3}
证明 \(\left[\mathbb Q(\sqrt2,\sqrt5):\mathbb Q\right]=4\)，并给出一组基。对每个 \(c\in\mathbb Q\)，求 \(\theta_c=\sqrt2+c\sqrt5\) 的最小多项式，并判断它何时生成整个 \(\mathbb Q(\sqrt2,\sqrt5)\)。
\end{exercise}
\begin{solution}
若 \(\sqrt5=a+b\sqrt2\)、\(a,b\in\mathbb Q\)，平方后比较 \(1,\sqrt2\) 的系数得 \(2ab=0\)、\(a^2+2b^2=5\)。\(b=0\) 将使 \(5\) 为有理平方，\(a=0\) 将使 \(5/2\) 为有理平方，二者均由素因数指数的奇偶性排除。所以 \(\sqrt5\notin\mathbb Q(\sqrt2)\)，二次塔的总次数为四，乘积基为 \(1,\sqrt2,\sqrt5,\sqrt{10}\)。

对任意 \(c\in\mathbb Q\)，令 \(\theta_c=\sqrt2+c\sqrt5\)。由 \(\theta_c^2=2+5c^2+2c\sqrt{10}\) 移项再平方，得 \(\bigl(\theta_c^2-(2+5c^2)\bigr)^2=40c^2\)。展开并整理常数项，得到首一关系
\[
P_c(\theta_c)=0,\qquad
P_c(X)=X^4-(4+10c^2)X^2+(2-5c^2)^2.
\]
当 \(c=0\) 时，\(P_0(X)=(X^2-2)^2\)，而 \(\theta_0=\sqrt2\) 的最小多项式是 \(X^2-2\)，它只生成二次子域。若 \(c\neq0\)，则 \(2-5c^2\neq0\)，因为 \(2/5\) 不是有理数的平方。由
\[
\theta_c(\sqrt2-c\sqrt5)=2-5c^2
\]
可知 \(\sqrt2-c\sqrt5=(2-5c^2)\theta_c^{-1}\in\mathbb Q(\theta_c)\)。将它与 \(\theta_c\) 相加、相减，再除以非零的 \(2\) 与 \(2c\)，便分别得到 \(\sqrt2\) 和 \(\sqrt5\)。故 \(\mathbb Q(\theta_c)=\mathbb Q(\sqrt2,\sqrt5)\)，其次数为四。由\cref{b1:prop:minimal-polynomial-degree}，\(\theta_c\) 的最小多项式次数也为四；它整除首一四次零化多项式 \(P_c\)，所以与 \(P_c\) 相等。这里确定最小多项式的依据，还包括 \(\theta_c\) 生成整个四次扩张。
\end{solution}

\begin{exercise}\label{b1:ex:14-4}
设 \(E/K\)、\(F/K\) 是同一扩张域中的有限扩张，次数互素。证明 \(E\cap F=K\)。反向命题是否成立？
\end{exercise}
\begin{solution}
由\cref{b1:thm:tower-law}，\([E\cap F:K]\) 同时整除 \([E:K]\)、\([F:K]\)，故为一。反向不成立：取 \(E=\mathbb Q(\sqrt2)\)、\(F=\mathbb Q(\sqrt3)\)。两者次数都为二；若交不为 \(\mathbb Q\)，二次塔公式迫使 \(E=F\)，但\cref{b1:prop:biquadratic-example}中 \(\sqrt3\notin E\) 的计算排除了这一点。所以交为 \(\mathbb Q\)，次数却不互素。

塔式公式给出了中间域次数必须满足的整除约束，但这些整数约束并不完全决定子域的位置关系。这里互素足以迫使交为基域，交为基域却不能反推出次数互素。
\end{solution}

\begin{exercise}\label{b1:ex:14-5}
设 \([L:K]=n<\infty\)，\(f\in K[x]\) 不可约且次数 \(d\) 与 \(n\) 互素。若 \(d>1\)，证明 \(f\) 在 \(L\) 中没有根；进一步判断：在上述次数互素的假设下，\(f\) 是否仍在 \(L[x]\) 中不可约？
\end{exercise}
\begin{solution}
若 \(a\in L\) 是根，则其最小多项式为 \(f\) 的首一化，故 \(\bigl[K(a):K\bigr]=d\)。塔式公式迫使 \(d\mid n\)，与互素且 \(d>1\) 矛盾。

判断 \(f\) 在 \(L[x]\) 中是否不可约，需要比较其根在 \(K\) 与 \(L\) 上的次数。关键是同一个大域可以按两条域塔计算次数：一条经过 \(L\)，另一条经过 \(K(a)\)。在包含 \(L\) 的扩张中取 \(f\) 的一个根 \(a\)，令 \(m=\bigl[L(a):L\bigr]\leq d\)。塔式公式给出
\[
mn=\bigl[L(a):K\bigr]=\bigl[L(a):K(a)\bigr]\,d.
\]
所以 \(d\mid mn\)，结合 \(\gcd(d,n)=1\) 得 \(d\mid m\)。于是 \(m=d\)，\(a\) 在 \(L\) 上的最小多项式仍有次数 \(d\)，只能等于 \(f\) 的首一化，故 \(f\) 仍不可约。这里添根可直接在 \(L[x]\) 中选 \(f\) 的一个不可约因子并取商域。
\end{solution}

\begin{exercise}\label{b1:ex:14-6}
设域 \(L\) 是两个子域 \(E,F\) 的并。证明 \(L=E\) 或 \(L=F\)。
\end{exercise}
\begin{solution}
若两者都是真子域，则既有 \(a\in E\setminus F\)，也有 \(b\in F\setminus E\)；否则一者包含另一者，并域就是后者。元素 \(a+b\) 属于 \(E\) 时，由减去 \(a\) 得 \(b\in E\)，矛盾；属于 \(F\) 时，由减去 \(b\) 得 \(a\in F\)，也矛盾。故不可能以两个真子域覆盖 \(L\)。
\end{solution}

\begin{exercise}\label{b1:ex:14-7}
设 \(K\subseteq E\subseteq L\)，且 \(E/K\) 为代数扩张。
\begin{enumerate}
\item 仅用代数性的传递性证明：对 \(a\in L\)，它在 \(E\) 上代数当且仅当在 \(K\) 上代数。
\item 利用超越次数的塔式公式证明 \(\operatorname{trdeg}_K L=\operatorname{trdeg}_E L\)。
\end{enumerate}
\end{exercise}
\begin{solution}
1.\ 在 \(K\) 上代数的方程也可视为 \(E\) 系数方程。反过来，若 \(a\) 在 \(E\) 上代数，则 \(E(a)/E\) 代数，结合 \(E/K\) 代数，由\cref{b1:thm:algebraic-tower}知 \(a\) 在 \(K\) 上代数。

2.\ \(E/K\) 为代数扩张，故 \(\operatorname{trdeg}_K E=0\)。将此代入\eqref{b1:eq:transcendence-tower}，得到 \(\operatorname{trdeg}_K L=\operatorname{trdeg}_E L\)。
\end{solution}

\begin{exercise}\label{b1:ex:14-8}
设 \(x,y\) 在 \(k\) 上代数独立。证明 \(x,x+y\) 也代数独立，并给出 \(k(x,y)\) 的另一组不是 \(k\)-线性坐标变换所得的超越基。
\end{exercise}
\begin{solution}
变量替换 \((X,Y)\mapsto(X,X+Y)\) 是多项式环的自同构，逆替换为 \((X,Y)\mapsto(X,Y-X)\)。所以非零多项式 \(P\) 代入 \((x,x+y)\) 后若为零，将产生 \(x,y\) 的非零关系，矛盾。另取 \(x,xy\)；它们生成的域包含 \(y=(xy)/x\)，故等于 \(k(x,y)\)。若二者相关，由 \(x\) 在 \(k\) 上超越及\cref{b1:lem:independence-adjoining}，\(xy\) 在 \(k(x)\) 上代数；于是 \(k(x,y)=k(x,xy)\) 在 \(k(x)\) 上代数，\cref{b1:lem:finite-transcendence-bound}便使该域的超越次数至多为一，与 \(x,y\) 独立矛盾。因此也是超越基。\(k\)-线性坐标变换所得的坐标必须形如 \(ax+by\)、\(a,b\in k\)；若 \(xy=ax+by\)，就会给出非零关系 \(XY-aX-bY=0\)。所以 \(xy\) 不是这样的线性式，所取超越基不由 \(k\)-线性坐标变换得到。
\end{solution}

\begin{exercise}\label{b1:ex:14-9}
设 \(x,y\) 在 \(k\) 上代数独立，令 \(s=x+y,t=xy\)。证明 \(s,t\) 代数独立，并证明 \(\bigl[k(x,y):k(s,t)\bigr]=2\)。
\end{exercise}
\begin{solution}
\(x,y\) 都是 \(X^2-sX+t\) 的根，且 \(y=s-x\)，故总次数至多为二。如果 \(s,t\) 相关，从两元素中可抽出的超越基至多有一个；\(k(x,y)/k(s,t)\) 代数于是迫使其超越次数至多为一，与 \(x,y\) 独立矛盾。因此 \(s,t\) 独立。

交换变量 \(x,y\) 给出 \(k(x,y)\) 的非恒等 \(k\)-自同构，它固定 \(s,t\) 而不固定 \(x\)。固定生成元的自同构会固定它们的一切有理表达式，因而逐点固定 \(k(s,t)\)。若 \(x\) 属于这个域，也必须被固定，矛盾。所以 \(x\notin k(s,t)\)，次数不能为一，故恰为二。这里用的是一个一般原则：属于被逐点固定的子域的元素，必被相应自同构固定。这一论证也适用于特征二。
\end{solution}

\begin{exercise}\label{b1:ex:14-10}
设 \(t\) 在 \(k\) 上超越，\(r\in k(t)\setminus k\)。证明 \(t\) 在 \(k(r)\) 上代数，进而证明 \(k(t)\) 中在 \(k\) 上代数的元素恰为 \(k\)。
\end{exercise}
\begin{solution}
写 \(r=A(t)/B(t)\)，其中 \(B\neq0\)。多项式 \(A(X)-r\cdot B(X)\in k(r)[X]\) 非零：若其系数全为零，取 \(B\) 的非零系数 \(b_i\) 就会得到 \(r=a_i/b_i\in k\)，矛盾。它以 \(t\) 为根，故 \(t\) 在 \(k(r)\) 上代数。若 \(r\) 还在 \(k\) 上代数，代数性的塔传递性会迫使 \(t\) 在 \(k\) 上代数，矛盾。所以一切非常数有理函数都超越，题目结论成立。

用\cref{b1:def:function-field}中的术语，这说明 \(k\) 在 \(k(t)\) 中相对代数闭；这个性质不要求 \(k\) 自身是代数闭域。
\end{solution}

\begin{exercise}\label{b1:ex:14-11}
设 \(t\) 超越，\(a,b,c,d\in k\) 且 \(ad-bc\neq0\)。证明 \(k(t)=k\bigl((at+b)/(ct+d)\bigr)\)，并说明当 \(ad-bc=0\) 且分母非零时发生什么。
\end{exercise}
\begin{solution}
令 \(u=(at+b)/(ct+d)\)。分母非零，因为 \(c,d\) 不同时为零且 \(t\) 超越。由 \((cu-a)t=b-du\) 得 \(t=(b-du)/(cu-a)\)。若 \(cu-a=0\)，则同时 \(b-du=0\)，推出 \(ad=bc\)，与假设矛盾，所以反解有效，两个域相等。若行列式为零且分母非零，两行 \((a,b),(c,d)\) 成比例，所给比值为 \(k\) 中常数，生成域退化为 \(k\)。
\end{solution}

\begin{exercise}\label{b1:ex:14-12}
设 \(T\) 是 \(k\) 上可数无限的代数独立集合，并将其不重复地枚举为 \(T=\{t_1,t_2,\ldots\}\)。证明 \(k(T)=\bigcup_{n\geq1}k(t_1,\ldots,t_n)\) 不是 \(k\) 上有限生成的扩张。
\end{exercise}
\begin{solution}
每个有理表达式只使用有限多个变量，故所给并等于 \(k(T)\)。若该域由有限多个元素生成，把这些元素的表达式所涉及变量合并，便有某个 \(N\) 使全部生成元都在 \(k(t_1,\ldots,t_N)\) 中，于是整个域都包含于该子域。但 \(t_{N+1}\) 在此子域上超越，特别地不在其中，矛盾。

关键是有限个域生成元总共只使用有限多个变量。无限超越基的基数证明也使用了同一有限支集思想：一个代数方程的有限多个系数，只涉及有限多个参数；这里则是有限多个有理表达式的变量范围可以同时限制在一个有限集合内。
\end{solution}

\begin{exercise}\label{b1:ex:14-13}
证明 \(A=k[X,Y]/(XY-1)\) 是整环，且其分式域同构于 \(k(t)\)。求该函数域在 \(k\) 上的超越次数。
\end{exercise}
\begin{solution}
映射 \(k[X,Y]\rightarrow k[t,t^{-1}]\)、\(X\mapsto t,Y\mapsto t^{-1}\) 满射。模 \(XY-1\)，每个混合单项式可以消去一个 \(XY\)，所以每个剩余类可写成 \(\sum_{i\geq0}a_iX^i+\sum_{j\geq1}b_jY^j\) 的有限和。其像为不同幂次的 Laurent 多项式，等于零当且仅当全部系数为零。因此核恰为 \((XY-1)\)，商环同构于 \(k[t,t^{-1}]\)，是整环；其分式域为 \(k(t)\)，超越次数为一。
\end{solution}

\begin{exercise}\label{b1:ex:14-14}
设整环 \(A\) 包含域 \(k\)，并由 \(k\) 与有限多个元素作为环生成，\(S\subseteq A\setminus\{0\}\) 为乘法集。证明 \(A\) 与 \(S^{-1}A\) 的分式域自然同构。由此说明只允许更多分母，不改变这里的超越次数。
\end{exercise}
\begin{solution}
由于 \(A\) 是整环，可把 \(S^{-1}A\) 视为 \(\operatorname{Frac}(A)\) 中由指定分式组成的子环。其分式域因此包含于 \(\operatorname{Frac}(A)\)。反方向，\(A\subseteq S^{-1}A\)，所以 \(A\) 的每个非零元素在 \(\operatorname{Frac}(S^{-1}A)\) 中可逆，\(A\) 的全部分式都属于后者。两个域由此相同，作为 \(k\) 扩张的超越次数当然相同。这个结论比较的是分式域，局部化前后的环未必同构。

分式域相同的证明只用了 \(A\) 是整环和 \(S\) 不含零。\(A\) 在 \(k\) 上有限生成的假设，保证其分式域是本节意义下的函数域，并未用于分式域相等本身。
\end{solution}
